\section{Preposiciones de tiempo (\textit{in, om, op})}
\textbf{Usos rápidos.}
\begin{itemize}[leftmargin=*]
  \item \textit{om} + hora exacta: \textit{om drie uur}. 
  \item \textit{op} + días/fechas: \textit{op maandag}, \textit{op 5 mei}. 
  \item \textit{in} + meses/estaciones/partes del día: \textit{in juli}, \textit{in de winter}, \textit{in de ochtend}. 
\end{itemize}

\textbf{Ejemplos}
\begin{itemize}[leftmargin=*]
  \item \textit{De les begint om acht uur}. 
  \item \textit{We zien elkaar op vrijdag}. 
  \item \textit{Ik ga op 10 juni op vakantie}. 
\end{itemize}

\textbf{Mini práctica}
\begin{itemize}[leftmargin=*]
  \item Escribe tres frases: una con \textit{om}, otra con \textit{op} y otra con \textit{in}.
  \item Pasa de fecha numérica a frase completa con preposición.
\end{itemize}
